\vspace{-2mm}
\section{Introduction}
\label{sec:introduction}
Large Language Models (LLMs) have shown impressive capabilities
in understanding, processing, and generating data, 
driving a wide range of data processing tasks, including data extraction, 
transformation, cleaning, and enrichment, as well as question-answering~\cite{zhao2024chat2data, huang2024transform, naeem2024retclean, buss2023generating, zhu2024autotqa} over structured~\cite{jin2022survey,nan2022fetaqa} and unstructured data~\cite{zhuang2023toolqa}. 
\rone{Recent declarative systems~\cite{shankar2024docetl,lin2024towards,patel2024lotus,liu2024declarative,wang2025unify,jo2024thalamusdb,satriani2025logical,sun2025quest,patel2024semantic,russo2025abacus}}, 
as well as SQL extensions to
data warehouses~\cite{databricks-llm,snowflake-llm, duckdb-llm,alloydb-llm} 
enable LLM-powered operators over textual columns, 
supporting similar tasks such as extraction, cleaning, and synthesis.
Most such LLM-powered data processing tasks or operators 
boil down to using an LLM 
to synthesize an answer for a given question 
from source data provided in-context, 
which could include 
structured or unstructured data. 
For example, data extraction can be framed as 
a natural language question, with the source data (documents or tables) provided in context. Or,  
in tabular question answering (TableQA)~\cite{jin2022survey,nan2022fetaqa}, 
the question is placed in the prompt, with
the relevant tables provided in context. 
In the following,
we refer to this context or source data as {\em text}, 
since both structured and unstructured data are converted 
into a textual form for use as context, comprising
one or more {\em sentences},
and to the data processing task or operator, expressed in natural language, as a {\em question}. 

% Large Language Models (LLMs) have demonstrated remarkable capabilities in understanding and processing textual data. Users or applications can now pose natural language questions with context in the form of textual passages or documents as input to an LLM, with the LLM generating answers. This question-answering capability powers many data-driven applications in database community. Data extraction and enrichment~\cite{zhao2024chat2data,naeem2024retclean} can be easily achieved by describing the task directly in the prompt. Natural language questions can now be answered over tables (known as TableQA~\cite{jin2022survey,nan2022fetaqa}) without SQL queries, democratizing users' access to the database. Database tuning~\cite{lao2025gptuner,zhou2023llm} uses LLMs to analyze DBMS characteristics, workloads, and historical query run logs, encoded as part of the context, to configure knobs, optimize query plans, and select indexes.  
% LLMs also drive data analytical systems over unstructured text~\cite{shankar2024docetl,lin2024towards,patel2024lotus,liu2024declarative} by evaluating declarative queries.  


When using LLMs for data processing, 
 it is {\em not clear how the LLM
response may have been derived, and if
the response can be trusted}.
Although LLM capabilities continue to improve, they are far from perfect~\cite{zhang2024benchmarking,liu2024lost}, often with accuracy less than 75\% in domains such as law, medicine, and finance~\cite{liang2023holistic, guha2023legalbench, islam2023financebench, arora2025healthbench, liu2024large}. %\agp{add better citations} 
%While accuracy or similar metrics are often used to evaluate the performance of LLMs on benchmarks, these are aggregate numbers and may not be reliable for a single invocation of the LLM. 
Gaining human trust in how answers
are produced and whether those answers are correct can be critical, 
especially when a wrong answer may lead to serious consequences. 
Consider a real-world use case: 
analysis of police use-of-force records, 
\anonymous{provided by our collaborator X\footnote{Name omitted for anonymity.}}{from the Police Records Access project, a collaboration that we are part of, along with various journalism and legal  partners~\cite{clean} to build the first-ever statewide
use-of-force database\footnote{Available at: \url{https://clean.sfchronicle.com} or \url{https://clean.latimes.com}}}.  
As in Figure~\ref{fig:example}, journalists aim to extract
names of officers involved in the use of force, from a 35-page document, and repeat this for 
over 3500 cases in the project, spanning over 1.5M pages~\cite{clean}.    
The LLM \code{gpt-4o-mini} returns the answer {\em Officer Anderson, Maya}. 
% Blindly using this result 
% could lead to serious negative consequences if the extraction is incorrect. 
%such as a false accusation. 
Even if the LLM has high overall accuracy,
that does not suffice to determine correctness for this specific query, 
where errors are especially likely
with large context sizes ~\cite{liu2024lost,bai2023longbench}. 
To ensure correctness,
our journalism collaborators 
often manually review each text document (of dozens of pages).
Here, an engineer on the project
developed a custom keyword search approach to identify potentially relevant
text portions for the answer (e.g., mentions of ``Officer'', ``Captain''), 
reducing review time by 7$\times$, but introducing many false positives and negatives. 
Overall, 
as the text size increases, 
verifying LLM answers becomes increasingly challenging. 
Ensuring human trust in answers is critical across domains, including analysis over medical, legal, financial, scientific, governmental, and journalistic documents. 
Effective mechanisms for verifying answer correctness 
are key to ensuring  usability of LLM-powered applications.  






%\yiming{add results for other datasets}

% \begin{figure}[tb]
%     \centering
%     \includegraphics[width=0.6\linewidth]{figures/provenance.pdf}
%     \vspace{-0.5em}
%     \caption{\small Provenance on Text. $T$, $Q$, $L$, and $A$ denote the Text, Question, LLM model, and Answer, respectively. $T'$ represents a portion of $T$, and $A'$ is the answer to $Q$ produced by $L$ over $T'$. } 
%     %\vspace{-1em}
%     \label{fig:provenance}
% \end{figure}

  
  

\begin{figure}[tb]
    \centering
    \includegraphics[width=0.9\linewidth]{figures/police_example_2.pdf}
    \vspace{-1em}
    \caption{\small Verifiable Provenance in Police Use of Force Records. (Actual values in police records have been replaced for privacy.)}  
    \vspace{-1em}
    \label{fig:example}
\end{figure}



% \begin{figure}[tb]
%     \centering
%     \includegraphics[width=1\linewidth]{figures/vanilla_LLM.pdf}
%     \vspace{-2em}
%     \caption{\small Agreement Percentage of LLM-generated Provenance. } %and rag-based baseline. }} 
%     \vspace{-1em}
%     \label{tag:agreement}
% \end{figure}

In this paper, we revisit 
a traditional database concept in the context of verification of LLM-generated answers: {\em provenance},
which, for now, we will informally refer to as a text portion 
from which the answer may have been derived. 
In the use of force analysis in Figure~\ref{fig:example}, 
if supporting {\em evidence} is provided, 
such as a fragment of the document that mentions the use of force types and 
their associations with the officer’s name, 
users can understand how the answer may have been generated 
and better assess its correctness. 


%current research on provenance. 
\topic{Current approaches to infer provenance are not reliable}  
To infer the provenance of an answer, one could leverage 
%\delete{work on} 
\revised{retrieval approaches by} identifying the top-$k$ sentences in the given text 
with the highest relevance scores to the question~\cite{sankararaman2024provenance,gao2023enabling,cohen2024contextcite}. 
These scores can be computed post hoc,
using embedding similarities~\cite{gao2023enabling} or next-token probabilities~\cite{cohen2024contextcite}.
However, these approaches do not guarantee
that the provenance is sufficient to reproduce the answer. 
\rthree{Table~\ref{tab:agreement} shows the accuracy of one such retrieval method based on embedding similarity, i.e., via RAG~\cite{lewis2020retrieval},  across five workloads: \code{Qasper}~\cite{dasigi2021dataset}, \code{NL\_DEV}~\cite{kwiatkowski2019natural},  \code{HotpotQA}~\cite{yang2018hotpotqa},  \code{CUAD}~\cite{hendrycks2021cuad}, and \code{PubMedQA}~\cite{jin2019pubmedqa}, each with 500 question-document pairs.} We vary the size of retrieved context (i.e., the inferred provenance) between 1\%, 5\%, and 10\% of the full text. 
To assess whether the retrieved context can reproduce the answer, we ask the same question to the same LLM (gpt-4o) with this context, and report accuracy as the fraction of cases where the resulting answer agrees with the one generated from the full text.\footnote{If the answers aren't identical, we use LLM-as-a-judge~\cite{gu2024survey,zheng2023judging} to determine accuracy, which is known to be effective at identifying semantically equivalent answers.} Unfortunately, {\em even when the retrieved context includes 10\% of the full text, accuracy is below 0.8}, making it neither reliable nor effective for human review. 




% \begin{table}[bt]
% \small
% \centering
% \scalebox{1}{
% \begin{tabular}{c|c|c|c}
% \hline
% & \textbf{Qasper} & \textbf{NL\_DEV} & \textbf{HotpotQA}  \\ \hline
% \textbf{gpt-4o-mini} & 0.62 & 0.68 & 0.65  \\ 
% \textbf{gpt-4o} & 0.65 & 0.77 & 0.79  \\ 
% \textbf{gemini-2-flash} & 0.41 & 0.37 & 0.61 \\ \hline 
% \end{tabular}}
% \caption{\small Agreement Percentage of LLM-generated Provenance.}
% \label{tab:agreement}
% \vspace{-3em}
% \end{table}

% \begin{table}[bt]
% \small
% \color{blue}
% \centering
% \scalebox{1}{
% \begin{tabular}{l|l|c|c|c}
% \hline
% \textbf{Method} & \textbf{Configuration} & \textbf{Qasper} & \textbf{NL\_DEV} & \textbf{HotpotQA} \\ \hline
% \multirow{3}{*}{\textbf{Retrieval}} 
% & \textbf{RAG-1\%} & 19\% & 17\% & 12\% \\ 
% & \textbf{RAG-5\%} & 46\% & 41\% & 35\% \\ 
% & \textbf{RAG-10\%} & 78\% & 69\% & 55\% \\ \hline
% \multirow{3}{*}{\textbf{LLM}} 
% & \textbf{gpt-4o-mini} & 62\% & 68\% & 65\% \\ 
% & \textbf{gpt-4o} & 65\% & 77\% & 79\% \\ 
% & \textbf{gemini-2-flash} & 41\% & 37\% & 61\% \\ \hline
% \end{tabular}}
% \caption{\small \revised{Accuracy of LLM-generated and retrieval-based provenance.}}
% \label{tab:agreement} 
% \vspace{-3em}
% \end{table}

\begin{table}[bt]
\small
%\color{blue}
\centering
\scalebox{0.76}{
\begin{tabular}{l|l|c|c|c|c|c}
\hline
\textbf{Method} & \textbf{Configuration} & \textbf{Qasper} & \textbf{NL\_DEV} & \textbf{HotpotQA} & \textbf{\rthree{CUAD}} & \textbf{\rthree{PubMedQA}} \\ \hline
\multirow{3}{*}{\textbf{Retrieval}} 
& \textbf{RAG-1\%} & 0.19 & 0.17 & 0.12 & \rthree{0.23} & \rthree{0.09} \\ 
& \textbf{RAG-5\%} & 0.46 & 0.41 & 0.35 & \rthree{0.39} & \rthree{0.24} \\ 
& \textbf{RAG-10\%} & 0.78 & 0.69 & 0.55  & \rthree{0.62} & \rthree{0.49}\\ \hline
\multirow{3}{*}{\textbf{LLM}} 
& \textbf{gpt-4o-mini} & 0.62 & 0.68 & 0.65 & \rthree{0.59} & \rthree{0.71}\\ 
& \textbf{gpt-4o} & 0.65 & 0.77 & 0.79 & \rthree{0.71} & \rthree{0.76} \\ 
& \textbf{gemini-2-flash} & 0.41 & 0.37 & 0.61 & \rthree{0.55} & \rthree{0.68}\\ \hline
%BLIP & BLIP & 1 & 1 & 1 \\ \hline 
\end{tabular}}
\caption{\small \revised{Accuracy of retrieval-based and LLM-generated provenance.}}
\label{tab:agreement} 
\vspace{-3.5em}
\end{table}




% \begin{table}[bt]
% \small
% \centering
% \scalebox{1}{
% \begin{tabular}{c|c|c|c}
% \hline
% & \textbf{Qasper} & \textbf{NL\_DEV} & \textbf{HotpotQA}  \\ \hline

% \end{tabular}}
% \caption{\small Agreement Percentage of RAG-generated Provenance.}
% \label{tab:agreement}
% \vspace{-3em}
% \end{table}

% \begin{table}[bt]
% \small
% \centering
% \scalebox{1}{
% \begin{tabular}{c|c|c|c}
% \hline
% & \textbf{Qasper} & \textbf{NL\_DEV} & \textbf{HotpotQA}  \\ \hline
% \textbf{RAG-1\%} & 0.19 & 0.17 & 0.12  \\ 
% \textbf{RAG-5\%} & 0.46 & 0.41 & 0.35  \\ 
% \textbf{RAG-10\%} & 0.78 & 0.69 & 0.55 \\ \hline 
% \end{tabular}}
% \caption{\small Agreement Percentage of RAG-generated Provenance.}
% \label{tab:agreement}
% \vspace{-3em}
% \end{table}


Another method is to directly
ask the LLM to identify the provenance of its answer
from the text, 
by prompting it to provide the line numbers
of the relevant sentences (by prefixing each sentence  with line numbers), 
as listed \papertext{in \vldb{~\cite{blip}.}} \techreport{in \code{LLM-provenance-prompt} below.} 
%During the generation of provenance, the LLM may fabricate sentences that are not present in the source document, even when explicitly instructed to use only its content, resulting in unfaithful provenance (e.g., the LLM-generated provenance is displayed in red box in Figure~\ref{fig:example}). \agp{This does not matter: you can add line numbers to the source doc, and ask an LLM to generate line numbers as part of its answer instead, and it will do that faithfully. Most people know the state of the art for citation is to ask LLMs to generate line numbers when they want actual citations, so it doesn't make sense to talk about a less sophisticated approach.} 
%In Table~\ref{tab:agreement}, we measure the accuracy 
% of LLM-generated provenance across three workloads: \code{Paper}, \code{HotpotQA}, and \code{NL\_DEV}, each with 500 question-document pairs. 
% To assess whether the provenance can reproduce the 
%  answer, we ask the same question to the same LLM
% with just the LLM-generated provenance as context,
% and report the percentage of cases 
% where the resulting answer agrees 
% with the one generated from the full text.\footnote{If the answers aren't identical, we use LLM-as-a-judge~\cite{gu2024survey,zheng2023judging} to determine agreement, which is known to be effective at identifying semantically equivalent answers.}
\revised{Unfortunately, as shown in Table~\ref{tab:agreement},} \code{gpt-4o-mini}, \code{gpt-4o}, and \code{gemini-2-flash} return provenance that reproduces the answer in only 65\%, 74\%, and 53\% of the cases, respectively---showing that LLMs are unable to identify the provenance 
that can reproduce the same (or equivalent) answer. \revised{When a provenance fails to reproduce the answer (e.g., those generated by retrieval or LLMs), human verification defaults to reading the entire document, significantly increasing human effort. 
} 

\techreport{
\noindent\underline{\code{LLM-provenance-prompt}}: \textcolor{blue}{[Question]}, \textcolor{blue}{[Answer]}, and \textcolor{blue}{[Context]} are placeholders for question $Q$, answer $A$, and the full text $T$.  }

\techreport{
\mypython{
\small 
    LLM-Provenance-Prompt: Given the following question: \textcolor{blue}{[Question]}, the corresponding answers are: \textcolor{blue}{[Answer]}. Your task is to extract the set of sentences from the provided context that contribute to generating these answers. Identify the most relevant sentences that support the given answers. Do not add explanations. Only return a list of sentence IDs. Do not return any words. The context is as follows: \textcolor{blue}{[Context]} }
}

Finally, other approaches~\cite{lee2025llm} 
require access to the model to compute attribution scores
for tokens based on training data but rely on access to training data, which 
is unavailable for black-box LLMs (e.g., OpenAI models) and frequently withheld for open-source models (e.g., Llama3). 
They also don't apply when operating on user-provided text, which is not part of the training data. 

%  \begin{figure}[tb]
%     \centering
%     %\vspace{-1em}
%     \includegraphics[width=0.6\linewidth]{figures/diagram.pdf}
%     \vspace{-1em}
%     \caption{\small Verifiable Provenance.\agp{ideally combine this with Figure~\ref{fig:example}}}  
%     \vspace{-2em}
%     \label{fig:diagram}
% \end{figure}

\topic{Verifiable provenance}
We use insights from the previous experiment
to introduce the concept of {\em verifiable provenance} over text  
for LLM-powered data processing tasks.
Instead of arbitrarily selecting a set of sentences 
based on their relevance to the question and answer, 
we require the provenance to be {\em verifiable}, 
i.e., the provenance is a subsequence of the text (e.g., green box in Figure~\ref{fig:example}) 
that can {\em reproduce the same (or equivalent) 
answer as the one obtained from the full text when queried using the same LLM}. 
In the example in Figure~\ref{fig:example}, 
the provenance in the green box 
is a subsequence of the original text, 
and asking the same question on just this provenance with the same LLM \code{gpt-4o-mini} 
leads to the same answer, {\em Officer Anderson, Maya}. 
Instead, the results in Table~\ref{tab:agreement} 
show that roughly {\em half of the LLM-generated provenance is not verifiable}. 

\topic{Conciseness of verifiable provenance and its challenges}
To help users easily assess answer correctness, 
the verifiable provenance must additionally be {\em small-sized}. 
Otherwise, the original text could itself serve as a provenance, but reviewing this is difficult and time-consuming.  
Identifying a verifiable  provenance that is also concise is nontrivial. 
%\sout{as it requires enumerating every subsequence of the text
%for the LLM to verify whether it produces an equivalent answer
%to that derived from the full text. }
One naive approach would involve considering all subsets of $k \in [0, n]$ sentences 
of the text with $n$ sentences \revised{and verifying whether the enumerated subset produces an equivalent answer with the LLM, which would scale as $2^n$. }
Another approach would 
start from the given text as the initial provenance, and 
iteratively remove a sentence at a time  
if the remaining still produces the same answer 
when using the same LLM---until no further sentences can be removed.
However, this approach incurs significant cost and latency, 
as it requires a number of LLM calls 
proportional to the number of sentences in the text. 
This baseline, on the Qasper dataset  with 9,000 tokens on average, 
takes $>$2 hours, incurring  $>$100$\times$ the 
cost of using an LLM for answering the question on the original text.



\topic{Merely a \sys to identify verifiable, concise provenance}  
We introduce \sys\footnote{Short for {\bf B}olt-on {\bf L}LM ver{\bf I}fiable {\bf P}rovenance.},
{\em a provenance discovery framework comprising 
a family of intuitive strategies to efficiently infer small-sized verifiable provenance}. 
We define the notion of {\em minimality} of verifiable provenance, 
where removing any sentence from the provenance 
results in a different answer.  
We further define
the notions of {\em strong} and {\em weak monotonicity} to
characterize the types of tasks for which
minimal provenance
provides additional guarantees, while
empirically showing that those tasks
dominate our workloads---and contrast
these notions  to those in the provenance
literature~\cite{cheney2009provenance,buneman2001and,green2007provenance,glavic2021data}.

As part of \sys, 
we present a two-phase approach that is {\em guaranteed to find a minimal verifiable provenance}. 
The first phase aggressively returns a small-sized verifiable provenance 
that may not be minimal.  
In this phase, we explore four strategies, 
including ranking text portions
based on their relevance to the question and answer, 
as well as different enumeration orders. 
The second phase takes the provenance returned from the first phase and  guarantees a minimal provenance. This two-phase approach, though intuitive, exhaustively explores 10 strategy combinations, each producing a minimal provenance. 
Using a cost model that we introduce, 
we present a theoretical cost comparison, and propose an adaptive strategy  
that integrates their best features. 
We further employ KV-caching~\cite{cai2024pyramidkv,liu2024minicache,adnan2024keyformer} to optimize each individual strategy to minimize inference cost. 
Finally, we extend \sys to support the discovery of multiple minimal
provenances rather than just one, to support use cases
where multiple pieces of evidence are helpful. 
% As we will demonstrate later in Section~\ref{sec:k-provenance}, 
% in some cases, exploring all provenances is essential to verify 
% the correctness of an answer. 
% To this end, we propose a top-$k$ provenance inference strategy as part of \sys that {\em guarantees the return of $k$ distinct minimal provenances}, enabling users to gain the most confidence in the answer's correctness.  

%\agp{update above and below based on new theoretical results}





%As mentioned earlier, a  minimal provenance may still be invalid, and LLMs alone cannot determine its validity, as they are used to verify whether the provenance reproduces the same answer. In contrast, users are often capable of judging the validity of a provenance upon inspection. For example, when users see Provenance 3 in Figure~\ref{fig:example}-2 under a {\em Reference} section, they can recognize that it does not correspond to the correct paper. Therefore, instead of returning a single provenance, we return the top-$k$ minimal provenances. If at least one is valid, the returned provenances can help users assess the correctness of the answer. Otherwise, the provided provenance is likely to signal that the answer is incorrect, as no valid provenance is returned. To support this, we further propose an algorithm that efficiently returns $k$ distinct minimal provenances, ranked by the likelihood that they are valid. 


Overall, \sys can be applied in a bolt-on manner, 
is compatible with any LLM model, 
and offers a theoretical guarantee 
of minimal provenance at a low cost.  Our contributions include:
\begin{itemize}[leftmargin=*]
    \item We {\em \revised{are the first} to introduce the notion of verifiable provenance for LLM-powered data processing}, outline key properties, including verifiability, minimality, and monotonicity, and 
    characterize the family of tasks where verifiable provenance applies.  
    \item We conduct an exhaustive exploration of intuitive  but effective strategies to efficiently infer minimal provenance at low cost.   
    \item We present a theoretical analysis of correctness and cost of the strategies, and use it to develop an adaptive strategy. 
    \item We perform comprehensive evaluations across \rthree{seven} datasets, demonstrating that it is possible to infer minimal provenance with  {\bf guaranteed 100\% accuracy}, achieving \textbf{over 30\%} higher accuracy than the best-performing baseline, at a low cost—comparable to that of performing the original task. 
\end{itemize}
\noindent We define the notion of provenance and present an  overview of \sys in Section~\ref{sec:def}. We then present our approaches for finding a single minimal provenance and top-$k$ minimal provenances in Section~\ref{sec:1-provenance} and  Section~\ref{sec:k-provenance}, respectively. We  evaluate \sys in Section~\ref{sec:exp}. 

%provenance size, cost and latency 

%In particular, consider a natual language question $Q$ over a text data $T$ (e.g., documents) by using an LLM model $M$, which returns the answer $A$. 







 